%% PATRÓN · Imagen con anotaciones: una imagen (render, captura) con una zona resaltada en coordenadas
%% RELATIVAS a la imagen (fracción de ancho y alto: no se rompe si cambia el tamaño) y una flecha curva
%% hacia otro elemento con su etiqueta. Si la imagen aún no existe, \figura deja un hueco.
%% Original: charla01/pipe.tex (tráquea -> pipe periódico).
\begin{tikzpicture}[font=\small]
  \def\H{4.0cm}\def\W{4.0cm}
  \node[anchor=south west,inner sep=0pt] (L) at (0,0) {\figura{figs/geometria.png}{\W}{\H}{Render de la geometría}};
  % zona resaltada: de (42 %, 60 %) a (63 %, 100 %) de la imagen
  \path (L.south west) ++(0.42*\W,0.60*\H) coordinate (Z0) (L.south west) ++(0.63*\W,\H) coordinate (Z1);
  \draw[acento,line width=1.2pt,rounded corners=3pt,fill=acento,fill opacity=0.10] (Z0) rectangle (Z1);
  \node[anchor=north,align=center,font=\footnotesize] at ($(L.south)+(0,-0.12)$)
    {lo real\\ {\color{tenue}dato clave · dato clave}};
  \draw[-{Stealth[length=7pt]},tenue,line width=1pt] ($(Z1|-Z0)!0.5!(Z1)+(0.15,0)$) to[out=10,in=170]
    node[above=2pt,pos=0.55,font=\scriptsize,text=tenue,fill=fondo,inner sep=1pt]{se simplifica a} (6.6,3.0);
  \node[draw=tinta,line width=0.6pt,rounded corners=3pt,fill=serieA!12!fondo,minimum width=4.4cm,minimum height=1.6cm,
        anchor=west,align=center] at (6.7,3.0) {el modelo\\ {\color{tenue}\footnotesize con referencia}};
\end{tikzpicture}
